% =====================================================================
% Projektantrag (Aufgabe 1) - Formblatt DIN A4
% F.-o. E-Projekt Mechatronik, WS 2026/27, Prof. Dr.-Ing. K. Pietsch
% Thema 15: Kraft-Weg-Messplatz fuer Kraftsensoren aus leitendem PLA
%
% Abgabe: 29.10.2026, 23:55, Moodle, PDF, von allen unterschrieben + gescannt
%
% Nutzung:
%   - Alle Stellen mit \feld{...} durch eigenen Text ersetzen.
%   - Graue Leitfragen (\hinweis{...}) helfen beim Schreiben.
%   - Fuer die Abgabe: \hinweisefalse setzen (Zeile unten) -> Leitfragen weg.
%   - KI-Regel (Anlage I): Texte selbst schreiben, jede KI-Nutzung im
%     Prompt-Protokoll (Anhang) eintragen.
% =====================================================================
\documentclass[11pt,a4paper]{article}

\usepackage[ngerman]{babel}
\usepackage[T1]{fontenc}
\usepackage[utf8]{inputenc}
\usepackage{lmodern}
\usepackage[a4paper,top=2.2cm,bottom=2.2cm,left=2.2cm,right=2.2cm]{geometry}
\usepackage{tabularx}
\usepackage{array}
\usepackage{booktabs}
\usepackage[table]{xcolor}
\usepackage{fancyhdr}
\usepackage{enumitem}
\usepackage{lastpage}
\usepackage{needspace}

% ---------------------------------------------------------------------
\newif\ifhinweise
\hinweisetrue   % <- fuer die Abgabe auf \hinweisefalse setzen

\definecolor{bhtgrau}{gray}{0.92}
\definecolor{hinweisgrau}{gray}{0.45}

% Leitfrage / Hinweis (nur sichtbar, solange \hinweisetrue)
\newcommand{\hinweis}[1]{\ifhinweise{\color{hinweisgrau}\small\itshape #1\par}\fi}
% Platzhalter fuer eigenen Text
\newcommand{\feld}[1]{\leavevmode\textcolor{red!70!black}{[#1]}}

% Ueberschrift eines Formularfeldes (graue Leiste)
\newcounter{feldnr}
\newcommand{\formfeld}[1]{%
  \par\medskip\Needspace{6\baselineskip}\noindent
  \stepcounter{feldnr}%
  \colorbox{bhtgrau}{\parbox{\dimexpr\linewidth-2\fboxsep}{\textbf{\thefeldnr.\ #1}}}%
  \par\smallskip}

\renewcommand{\arraystretch}{1.3}
\setlength{\parindent}{0pt}
\setlist{nosep,leftmargin=1.2em}

\pagestyle{fancy}
\fancyhf{}
\lhead{\small Projektantrag | F.-o. E-Projekt Mechatronik}
\rhead{\small Berliner Hochschule für Technik}
\cfoot{\small Seite \thepage{} von \pageref{LastPage}}

% =====================================================================
\begin{document}

\begin{center}
  {\Large\bfseries Interner Projektantrag}\\[2pt]
  {\large Laboreinrichtung / Funktionsmuster eines mechatronischen Systems}\\[2pt]
  {\small Modul M08 F.-o. E-Projekt · Master Mechatronik · WS 2026/27}
\end{center}

% ---------------------------------------------------------------------
\formfeld{Zum Projekt}
\begin{tabularx}{\linewidth}{@{}>{\bfseries}p{4.2cm}X@{}}
Projekttitel     & Entwicklung eines Kraft-Weg-Messplatzes für Kraftsensoren aus leitendem PLA als „DMS“ (Thema 15) \\
Auftraggeber     & \feld{z.\,B. Berliner Hochschule für Technik, FB VII, vertreten durch \dots} \\
Betreuer         & Prof. Dr.-Ing. Karsten Pietsch \\
Labor            & \feld{z.\,B. B402L (GOS)} \\
Projektlaufzeit  & \feld{von \dots\ bis \dots} \\
Antragsdatum     & \feld{TT.MM.2026} \\
\end{tabularx}

% ---------------------------------------------------------------------
\formfeld{Projektteam (Kernteam)}
\hinweis{Personelle Voraussetzungen: Wer bringt welche Kompetenz mit (CAD, Elektronik, Programmierung, 3D-Druck)? Welche Teamrolle nach Belbin? Wer übernimmt welche Arbeitspakete (AP\,1--5)?}
\begin{tabularx}{\linewidth}{|l|l|X|l|}
\hline
\rowcolor{bhtgrau}\textbf{Name} & \textbf{Rolle im Projekt} & \textbf{Kompetenzen / Voraussetzungen} & \textbf{AP} \\ \hline
Khoa Nguyen            & \feld{\dots} & \feld{\dots} & \feld{\dots} \\ \hline
Joffre Nunez Rodriguez & \feld{\dots} & \feld{\dots} & \feld{\dots} \\ \hline
\end{tabularx}

% ---------------------------------------------------------------------
\formfeld{Ziel (nach SMART-Regel)}
\hinweis{Ein Gesamtziel in 1--2 Sätzen, danach die fünf Kriterien prüfen. Tipp: Messbar heißt mit Zahl und Einheit (Kraftbereich in N, Weg in mm, Auflösung, Wiederholgenauigkeit \dots).}
\textbf{Projektziel:} \feld{\dots}

\smallskip
\begin{tabularx}{\linewidth}{|>{\bfseries}l|X|}
\hline
\rowcolor{bhtgrau}\textbf{Kriterium} & \textbf{Ausprägung im Projekt} \\ \hline
Spezifisch   & \feld{Was genau wird gebaut?} \\ \hline
Messbar      & \feld{Woran sieht man, dass das Ziel erreicht ist? (Kennwerte)} \\ \hline
Attraktiv / Akzeptiert & \feld{Nutzen für Stakeholder: Laborkapazität, Demonstrator für die Lehre} \\ \hline
Realistisch  & \feld{Warum mit Team, Budget und Zeit machbar?} \\ \hline
Terminiert   & \feld{Bis wann? (Funktionsmuster, Abschluss)} \\ \hline
\end{tabularx}

% ---------------------------------------------------------------------
\formfeld{Aufgabenstellung}
\textbf{Technisches Gebilde (Teilaufgabe 1):}
\hinweis{Welches mechatronische System wurde zugeteilt? Was ist die Ausgangssituation im Labor (was existiert schon, was fehlt)?}
\feld{\dots}

\medskip
\textbf{Entwicklungsthema (Teilaufgabe 2):}
\hinweis{Mindestens drei nachvollziehbare Sätze. Alle Schnittstellen des mechatronischen Systems nennen: Mechanik (Einspannung, Führung, Messaufnehmer), Elektrotechnik (Sensoranbindung, Messschaltung, Mikrocontroller), Software (Steuer- und Auswertesoftware), Mensch-Maschine (Bedienung). Bezug: Elgeneidy et al.\ 2018.}
\feld{\dots}

\medskip
\textbf{Kurzbeschreibung des Vorhabens:}
\hinweis{Wie wollt ihr vorgehen? Kurzer Verweis auf die Arbeitspakete AP\,1--5 genügt (ggf. als Anlage beifügen).}
\feld{\dots}

% ---------------------------------------------------------------------
\formfeld{Zu erarbeitende Ergebnisse}
\hinweis{Liefergegenstände = was am Ende übergeben wird. Erfolgskriterium = woran man prüft, dass es gut ist (möglichst messbar).}
\begin{tabularx}{\linewidth}{|l|X|X|}
\hline
\rowcolor{bhtgrau}\textbf{Nr.} & \textbf{Liefergegenstand} & \textbf{Erfolgskriterium} \\ \hline
E1 & \feld{z.\,B. Funktionsmuster Messplatz} & \feld{\dots} \\ \hline
E2 & \feld{\dots} & \feld{\dots} \\ \hline
E3 & \feld{\dots} & \feld{\dots} \\ \hline
E4 & \feld{\dots} & \feld{\dots} \\ \hline
\end{tabularx}

% ---------------------------------------------------------------------
\formfeld{Kosten}
\hinweis{Gesamtbudget aus Teilbudgets mit ersten wirtschaftlichen Überlegungen. Sachkosten: Steckbrief nennt < 200\,€. Personalkosten: Stunden je Person × angenommener Stundensatz (Annahme begründen). Eine Reserve ist sinnvoll.}
\begin{tabularx}{\linewidth}{|X|l|r|}
\hline
\rowcolor{bhtgrau}\textbf{Position} & \textbf{Annahme / Grundlage} & \textbf{Betrag} \\ \hline
\multicolumn{3}{|l|}{\textit{Sachkosten}} \\ \hline
\feld{Mechanik} & \feld{\dots} & \feld{\dots\ €} \\ \hline
\feld{Elektronik / Sensorik} & \feld{\dots} & \feld{\dots\ €} \\ \hline
\feld{Verbrauchsmaterial (Filament \dots)} & \feld{\dots} & \feld{\dots\ €} \\ \hline
\feld{Reserve} & \feld{\dots\,\%} & \feld{\dots\ €} \\ \hline
\multicolumn{2}{|r|}{\textbf{Summe Sachkosten}} & \feld{\dots\ €} \\ \hline
\multicolumn{3}{|l|}{\textit{Personalkosten}} \\ \hline
\feld{Projektteam} & \feld{\dots\ h × \dots\ €/h} & \feld{\dots\ €} \\ \hline
\multicolumn{2}{|r|}{\textbf{Gesamtbudget}} & \feld{\dots\ €} \\ \hline
\end{tabularx}

% ---------------------------------------------------------------------
\formfeld{Randbedingungen}
\hinweis{Grenzen des Projekts: Budget, Zeit (ein Semester), Labor und Öffnungszeiten, verfügbare Geräte (3D-Drucker, Messtechnik), Bestellschluss 05.11.2026, Sicherheit, was ausdrücklich \emph{nicht} Teil des Projekts ist.}
\begin{itemize}
  \item \feld{\dots}
  \item \feld{\dots}
  \item \feld{\dots}
\end{itemize}

% ---------------------------------------------------------------------
\formfeld{Termine / Meilensteine}
\hinweis{Mindestens vier, nicht nur Semestertermine: eigene Meilensteine ergänzen (z.\,B. Sensor kontaktiert, Messschaltung läuft, Kalibrierung abgeschlossen). Die Semestertermine stehen in Semester\_Aktuelle\_Planung.pdf, S.\,6.}
\begin{tabularx}{\linewidth}{|l|l|X|}
\hline
\rowcolor{bhtgrau}\textbf{Nr.} & \textbf{Termin} & \textbf{Meilenstein / Ergebnis} \\ \hline
M1 & \feld{TT.MM.} & \feld{\dots} \\ \hline
M2 & \feld{TT.MM.} & \feld{\dots} \\ \hline
M3 & \feld{TT.MM.} & \feld{\dots} \\ \hline
M4 & \feld{TT.MM.} & \feld{\dots} \\ \hline
M5 & \feld{TT.MM.} & \feld{\dots} \\ \hline
M6 & \feld{TT.MM.} & \feld{\dots} \\ \hline
\end{tabularx}

% ---------------------------------------------------------------------
\formfeld{Risikoanalyse}
\hinweis{Kurz und sinnvoll: technische, terminliche, finanzielle und personelle Risiken. Wahrscheinlichkeit und Auswirkung z.\,B. mit gering / mittel / hoch bewerten, dazu eine Gegenmaßnahme.}
\begin{tabularx}{\linewidth}{|l|X|c|c|X|}
\hline
\rowcolor{bhtgrau}\textbf{Nr.} & \textbf{Risiko} & \textbf{W.} & \textbf{A.} & \textbf{Gegenmaßnahme} \\ \hline
R1 & \feld{\dots} & \feld{\dots} & \feld{\dots} & \feld{\dots} \\ \hline
R2 & \feld{\dots} & \feld{\dots} & \feld{\dots} & \feld{\dots} \\ \hline
R3 & \feld{\dots} & \feld{\dots} & \feld{\dots} & \feld{\dots} \\ \hline
R4 & \feld{\dots} & \feld{\dots} & \feld{\dots} & \feld{\dots} \\ \hline
\end{tabularx}

\smallskip
{\small W.\ = Eintrittswahrscheinlichkeit, A.\ = Auswirkung}

% ---------------------------------------------------------------------
\formfeld{Unterschriften}
Mit ihrer Unterschrift bestätigen alle Mitglieder des Projektteams den Inhalt dieses Projektantrags.

\vspace{1.2cm}
\begin{tabularx}{\linewidth}{XX}
\rule{0.9\linewidth}{0.4pt} & \rule{0.9\linewidth}{0.4pt} \\
Ort, Datum, Khoa Nguyen & Ort, Datum, Joffre Nunez Rodriguez \\
\end{tabularx}

% =====================================================================
\newpage
\section*{Anhang: Erklärung zur Verwendung von KI}

Für die Vorbereitung dieses Projektantrages wurden KI-Werkzeuge eingesetzt. Die Verantwortung für die finalen Inhalte liegt vollständig bei den Autoren.

\subsection*{Prompt-Protokoll}
\hinweis{Jede KI-Nutzung über Rechtschreib- und Grammatikprüfung hinaus eintragen. Die erste Zeile dokumentiert die Nutzung zur Erstellung dieser Vorlage.}
\begin{tabularx}{\linewidth}{|p{2.6cm}|l|X|X|}
\hline
\rowcolor{bhtgrau}\textbf{Tool / Modell} & \textbf{Datum} & \textbf{Einsatzzweck} & \textbf{Prompt (wesentlicher Teil)} \\ \hline
Claude (Anthropic) & 15.10.2026 & Zusammenfassung der Aufgabenstellung, Gliederung des Formblatts mit Leitfragen (ohne Inhaltstexte) & „lo hago con otro compañero y necesito hacerlo según este proyecto [Steckbrief Thema 15]“ \\ \hline
\feld{\dots} & \feld{\dots} & \feld{\dots} & \feld{\dots} \\ \hline
\feld{\dots} & \feld{\dots} & \feld{\dots} & \feld{\dots} \\ \hline
\end{tabularx}

\end{document}